%
%
%
%
%
%
%
%
%
%
%
%
%
%
%
%
%
%
%
%
%
%
%
%
%
%
%
%
%
%
%

% ── Platz: Grundgrösse, Dichte, Zeilenhöhe, Umbruchreserve ───────────
% Müssen VOR styles/30_layout_spacing stehen.
%
% Drei Hebel, drei verschiedene Fragen — bewusst nicht einer:
%   Grundgrösse   wie gross die Schrift ist
%   Dichte        wie viel Leerraum zwischen den Blöcken liegt
%   Zeilenhöhe    wie eng die Zeilen aufeinander folgen
% Sie sind getrennt, weil unterschiedlich riskant: Abstände sind Leerraum und
% vertragen jede Skalierung, Zeilenhöhe enthält die Schrift selbst — nur in
% kleinen Schritten senken und danach das PDF auf kollidierende Formelzeilen
% prüfen.
%
% Die Grundgrösse des Satzes. Leer = die Klassenoption aus main.tex
% (\documentclass[8pt]). Nimmt den ganzen Schriftsatz mit — Klassenstufen,
% Titel- und Balkengrössen, Diagramm-Beschriftungen —, sodass die Verhältnisse
% erhalten bleiben. Erreicht auch Grössen, die extsizes nicht als Klassenoption
% anbietet (7pt). Für eine ZSF, die den Platz nicht über Dichte und Zeilenhöhe
% allein bekommt: der grösste der drei Hebel, ~30 % Fläche zwischen 8pt und 7pt.
% \def\ZSFBodyFontSize{7pt}       % leer = Klassengrösse (8pt)
% \def\ZSFDensityFactor{0.85}     % 1    = Referenz, kleiner = dichter
%
% 0.90 ist der Wert, mit dem diese ZSF vor der Template-Migration gesetzt war
% (damals als rohes \linespread{0.9} in styles/00_packages). Die Referenz 0.95
% brachte rund 5 % mehr Zeilenhöhe über den gesamten Satz — in den ~50 mm
% schmalen Spalten der sichtbarste Teil des Seitenwachstums. Hier statt dort:
% Zeilenhöhe ist Layout und keine Paket-Einstellung.
\def\ZSFLeadingFactor{0.90}       % 0.95 = Referenz, kleiner = enger
%
% Die zwei ROLLEN-Faktoren auf die Grundgrösse. Sie beantworten dieselbe Frage
% für die zwei Inhaltsarten — wie laut steht dieser Inhalt neben den
% Titelbalken, die ihn gliedern? — und sind deshalb getrennt von
% \ZSFBodyFontSize, das die Balken mitnähme.
%
% Boxinhalt: Boxtext, Tabellenzelle, Formel. 0.875 = 7pt bei 8pt Grundgrösse,
% und das ist die Grösse, in der diese ZSF vor der Template-Migration gesetzt
% war. Sie ist hier keine Geschmacksfrage, sondern eine Passform: Eine
% Formelzelle einer zweispaltigen Tabelle ist 86.2 pt breit, und
% sin(arctan(x)) = x/sqrt(1+x^2) braucht bei 8pt Inhalt 86.3 pt — es bricht um.
% Bei 7pt sind es 75.9 pt. Über \ZSFBodyFontSize gelöst schrumpften die
% Titelbalken von 8.1 auf 7.1 bp mit; die stehen aber richtig.
\def\ZSFContentScale{0.875}       % 1 = Klassengrösse (8pt)
%
% Fliesstext: `runintext`, freier Text zwischen den Bausteinen, das Register.
% Gleicher Wert wie der Boxinhalt = beide auf 7pt, der Zustand vor der
% Migration. Wer die Prosa LEISER als den Boxinhalt will, senkt nur diesen
% Faktor — die Box trägt den prüfungsrelevanten Inhalt, der Fliesstext
% verbindet.
\def\ZSFFlowTextScale{0.875}      % 1 = Klassengrösse (8pt)
%
% Bereichsfaktoren multiplizieren den globalen Faktor für IHREN Bereich —
% für ZSF, die ungleich verteilt sind (box-lastig vs. textlastig).
\def\ZSFDensityBoxes{0.85}        % Innen-/Aussenabstände der Inhaltsboxen
% \def\ZSFDensityText{1}          % Absatzabstand im Fliesstext
% \def\ZSFDensityTables{1}        % Zell- und Zeilenabstand in Tabellen
% \def\ZSFDensityStructure{1}     % Polsterung und Abstände der Titelbalken
%
% Umbruchsteuerung an Titelbalken: in dieser ZSF komplett aus, siehe
% \ZSFBreakReservesOff nach den \input-Zeilen unten. Der Faktor bleibt
% deshalb auf der Referenz — er hat hier nichts mehr zu skalieren.
% \def\ZSFBreakReserveFactor{1}   % 1 = Referenz
%
% Der Abstand NACH einem Titelbalken — der ganze Abstand bis zum ersten
% Block darunter, Fliesstext wie Box. Eigener Regler und nicht Teil von
% \ZSFDensityStructure: Wie eng eine Überschrift an ihrem Inhalt klebt, ist
% eine Bindungs-Entscheidung und keine Frage der Gesamtdichte — in einer
% ZSF mit vielen kurzen Abschnitten will man ihn enger als alles andere.
\def\ZSFBarGapFactor{0.5}         % 1 = Referenz (3pt), kleiner = enger
%
% Schriftgrösse der Diagramm-Beschriftung — beide Stufen gemeinsam
% (\ZSFfontDiagramLabel und \ZSFfontDiagramLabelSmall). Kleiner für Fächer mit
% vielen gedrängten Zeichnungen; die zwei Stufen bleiben dabei unterscheidbar.
% Ein ROLLEN-Faktor: Er multipliziert die Grundgrösse für diese eine Rolle,
% statt sie zu ersetzen — wie die Bereichsfaktoren bei der Dichte.
% \def\ZSFDiagramLabelScale{0.85} % 1 = Referenz, kleiner = kleinere Labels

% ============================================================
% PACKAGES
% ============================================================

\usepackage[T1]{fontenc}
\usepackage[utf8]{inputenc}
\usepackage[ngerman]{babel}
\usepackage[sfdefault,lf]{carlito} % Carlito ist in LaTeX das Äquivalent zu Calibri (Open-Source von Google), 'lf' erzwingt Lining Figures (normale Zahlenhöhe) statt Old-Style (Zahlen wie ein kleines o)
\usepackage{sansmath}          % Mathe-Modus auch in Sans-Serif
\sansmath
\linespread{0.9} % Etwas engere Zeilen für mehr Kompaktheit bei 10pt

\usepackage[landscape, a4paper, margin=1.5mm]{geometry}
\usepackage{microtype}
\usepackage{multicol}

\usepackage{mathtools}          % Superset von amsmath: DeclarePairedDelimiter, coloneqq, etc.
\let\Bbbk\relax % Verhindert Konflikt zwischen newpxmath und amssymb
\usepackage{amssymb}
\usepackage{siunitx}
\sisetup{
  mode=match,
  propagate-math-font=true,
  reset-text-family=false,
  exponent-product = \cdot,
  output-decimal-marker = {,},
  per-mode = symbol
}

\usepackage[table]{xcolor}
\usepackage{array}
\usepackage{tabularx}
\usepackage{tabularray}
\usepackage[column=O]{cellspace}
\usepackage{ragged2e} % Erlaubt Silbentrennung bei linksbündigem/zentriertem Flattersatz in schmalen Spalten
\usepackage{xstring}
\usepackage{varwidth}

\usepackage{graphicx}
\usepackage{tikz}
\usepackage{pgfplots}
\pgfplotsset{compat=1.18}
\usepackage{tcolorbox}
\tcbuselibrary{skins, breakable, raster}

\usepackage{hyperref}
\usepackage{enumitem}
\usepackage{makeidx}               % Stichwortverzeichnis (Setup in styles/65_index)

% ============================================================
% MATH MACROS
% ============================================================

% --- Zahlmengen (Blackboard Bold) ---
\newcommand{\R}{\mathbb{R}}
\newcommand{\C}{\mathbb{C}}
\newcommand{\N}{\mathbb{N}}
\newcommand{\Z}{\mathbb{Z}}
\newcommand{\Q}{\mathbb{Q}}

% --- Differentialoperator (aufrechtes d) ---
\newcommand{\dd}{\,\mathrm{d}}

% --- Operatornamen ---
\DeclareMathOperator{\sgn}{sgn}
\DeclareMathOperator{\grad}{grad}
\DeclareMathOperator{\divg}{div}
\DeclareMathOperator{\rot}{rot}
\newcommand{\vect}[1]{\vec{#1}}

% --- Gepaarte Begrenzer (mathtools) ---
\DeclarePairedDelimiter{\abs}{\lvert}{\rvert}
\DeclarePairedDelimiter{\norm}{\lVert}{\rVert}

% --- Pfeile (Fix für sansmath Überlappungen) ---
% Da sansmath/Carlito bei zusammengesetzten langen Pfeilen (wie \Longrightarrow) 
% unschöne Überlappungen erzeugt, biegen wir diese auf die sauberen kurzen Pfeile um.
% Das spart auf dem Cheat Sheet zudem Platz.
\renewcommand{\implies}{\;\Rightarrow\;}
\renewcommand{\iff}{\;\Leftrightarrow\;}
\renewcommand{\Longrightarrow}{\Rightarrow}
\renewcommand{\Longleftrightarrow}{\Leftrightarrow}
\renewcommand{\longrightarrow}{\rightarrow}
\renewcommand{\longleftarrow}{\leftarrow}

% --- Summen/Limes-Mode-System ---
\newcommand{\ZSFsumMode}{auto}
\newcommand{\SetZSFsumMode}[1]{\renewcommand{\ZSFsumMode}{#1}}
\newcommand{\ZSFsumSide}[2]{\sum\nolimits_{#1}^{#2}}
\newcommand{\ZSFsumStack}[2]{\sum\limits_{#1}^{#2}}
\newcommand{\ZSFsumAuto}[2]{%
  \IfStrEqCase{\ZSFsumMode}{%
    {side}{\ZSFsumSide{#1}{#2}}%
    {stack}{\ZSFsumStack{#1}{#2}}%
    {auto}{\sum_{#1}^{#2}}%
  }[\sum_{#1}^{#2}]%
}

\newcommand{\ZSFlimMode}{auto}
\newcommand{\SetZSFlimMode}[1]{\renewcommand{\ZSFlimMode}{#1}}
\newcommand{\ZSFlimAuto}[1]{%
  \IfStrEqCase{\ZSFlimMode}{%
    {side}{\lim\nolimits_{#1}}%
    {stack}{\lim\limits_{#1}}%
    {auto}{\lim_{#1}}%
  }[\lim_{#1}]%
}

% --- Semantisches Formel-Highlighting ---
% Verknuepft visuell zusammengehoerige Teile einer Formel ueber mehrere
% Umformungsschritte hinweg. Farben zentral in 40_colors_structure.tex.
\newcommand{\ZSFmhlA}[1]{{\color{MathHLPrimary}#1}}
\newcommand{\ZSFmhlB}[1]{{\color{MathHLSecondary}#1}}
\newcommand{\ZSFmhlC}[1]{{\color{MathHLTarget}#1}}
\newcommand{\ZSFmhlD}[1]{{\color{MathHLTertiary}#1}}


% ── Formelsatz: Summen- und Limes-Grenzen ────────────────────────────
% side = neben dem Zeichen (spart Höhe), stack = darüber/darunter,
% auto = LaTeX entscheidet nach Kontext.
% \SetZSFsumMode{auto}            % auto | side | stack
% \SetZSFlimMode{auto}            % auto | side | stack

% ============================================================
% 11_math_advanced.tex — OPTIONALES Math-Erweiterungsmodul
% ============================================================
% Opt-in: in preamble.tex einkommentieren für mathe-lastige
% Fächer (LinAlg, Analysis). Baut auf 10_math.tex auf und
% ergänzt nur das Delta — keine Basis-Makros duplizieren.
%
% Voraussetzungen (aus 00_packages.tex):
%   - mathtools, unicode-math   (Operatoren, Symbole)
%   - tikz                      (drawbrace / tikzmark / annote)
% ============================================================

% --- Lineare-Algebra-Operatoren ---
\DeclareMathOperator{\Ker}{Ker}
\DeclareMathOperator{\rang}{rang}
\DeclareMathOperator{\Spur}{Spur}
\DeclareMathOperator{\tr}{tr}
\DeclareMathOperator{\diag}{diag}
\DeclareMathOperator{\spanop}{span}
\DeclareMathOperator{\eig}{eig}
\DeclareMathOperator{\algVfh}{algVfh}
\DeclareMathOperator{\gVfh}{gVfh}
\DeclareMathOperator{\proj}{proj}
\DeclareMathOperator{\sig}{sig}
\DeclareMathOperator{\Adj}{Adj}

% --- Analysis-Operatoren (Abweichung dieses Forks) ---
% Grad(p) ist der Polynomgrad und bewusst gross geschrieben — \grad aus
% 10_math ist der Gradient und eine andere Sache.
% Die Areafunktionen sind die Umkehrungen der Hyperbelfunktionen; TeX bringt
% dafuer nichts mit, und \operatorname im Kapitel meldet der Linter zu Recht:
% Ein Operator, der mehrfach vorkommt, gehoert einmal hierher.
\DeclareMathOperator{\Grad}{Grad}
\DeclareMathOperator{\Arsinh}{Arsinh}
\DeclareMathOperator{\Arcosh}{Arcosh}
\DeclareMathOperator{\Artanh}{Artanh}

% --- Phasenportraet / GGWP Semantik-Typen (Ink-konform) ---
\newcommand{\ZSFggwpKnoten}[1][Knoten]{\ZSFInk{ChapterColor9}{\textbf{#1}}}
\newcommand{\ZSFggwpSattel}[1][Sattel]{\ZSFInk{ChapterColor14}{\textbf{#1}}}
\newcommand{\ZSFggwpStrudel}[1][Strudel]{\ZSFInk{ChapterColor4}{\textbf{#1}}}
\newcommand{\ZSFggwpZentrum}[1][Zentrum]{\textbf{#1}}

% --- Gotisches Im/Re -> moderne aufrechte Operatoren ---
% AtBeginDocument, weil unicode-math seine Symbole ebenfalls erst dort
% definiert — eine Praeambel-Renewcommand wuerde still ueberschrieben.
\AtBeginDocument{%
  \let\oldIm\Im
  \renewcommand{\Im}{\operatorname{Im}}%
  \let\oldRe\Re
  \renewcommand{\Re}{\operatorname{Re}}%
}

% --- TikZ-Klammern für Matrix-Annotationen ---
\usetikzlibrary{decorations.pathreplacing,calc}
% Das optionale Argument sind KNOTEN-OPTIONEN (z.B. xshift), keine Länge. Der
% Default ist deshalb leer: Ein Mass an dieser Stelle liest pgf als
% Pfeilspezifikation und bricht ab ("Unknown arrow tip kind"). Der vertikale
% Versatz der Klammer hängt ohnehin an den Koordinaten, siehe unten.
\newcommand{\tikzmark}[2][]{\tikz[remember picture, overlay, baseline=-0.5ex]\node[#1](#2){};}
% Die Marker von \tikzmark sitzen auf der GRUNDLINIE des Terms. Ohne Versatz
% liegt die Klammer-Sehne genau dort und läuft als waagrechte Linie mitten
% durch die Formel. Der Versatz muss an den Koordinaten hängen, nicht als
% yshift an der Pfadoption: Knoten aus einem anderen Bild (remember picture)
% werden absolut aufgelöst und ignorieren die Pfad-Transformation.
%
% Richtung und Dekoration gehören zusammen, deshalb setzt jeder Stil beides.
% \ZSF@braceSign trägt das Vorzeichen zur Koordinate — die Optionen werden vor
% den Koordinaten ausgewertet, der Wert steht also rechtzeitig.
\makeatletter
\newcommand{\ZSF@braceSign}{-}
\tikzset{
  brace/.style={decorate, decoration={brace},
                /utils/exec=\renewcommand{\ZSF@braceSign}{}},
  brace mirrored/.style={decorate, decoration={brace,mirror},
                /utils/exec=\renewcommand{\ZSF@braceSign}{-}},
}
\newcounter{brace}
\setcounter{brace}{0}
% \drawbrace[stil]{marke links}{marke rechts}
%   brace           — Klammer ÜBER dem Term (Default)
%   brace mirrored  — Klammer UNTER dem Term
\newcommand{\drawbrace}[3][brace]{%
  \refstepcounter{brace}%
  \tikz[remember picture, overlay]\draw[#1]
    ([yshift=\ZSF@braceSign\ZSFbraceDrop]#2.center)--%
    ([yshift=\ZSF@braceSign\ZSFbraceDrop]#3.center)%
    node[pos=0.5, name=brace-\thebrace]{};%
}
\makeatother
\newcommand{\annote}[3][]{%
  \tikz[remember picture, overlay]\node[#1] at (#2) {#3};%
}

%
%
%
%
%
%
%
%
%
%
\newcommand{\ZSFbraceRoom}{\vspace{\dimexpr\ZSFbraceDrop+\ZSFspaceM\relax}}

% --- Asterisk-Definitionen (mathabx-Namen auf Standard-Symbole) ---
\newcommand{\asterisk}{\ast}
\newcommand{\oasterisk}{\circledast}
   % opt-in: aufrechtes Im/Re, Grad, Area-Funktionen
\input{styles/12_plots}           % opt-in: pgfplots (Funktionenübersicht Anhang)
% ============================================================
% TABLE COLUMNS AND TABLE HELPERS
% ============================================================

\renewcommand{\tabularxcolumn}[1]{m{#1}} % Zentriert X-Spalten vertikal für Formeln
\newcolumntype{C}{>{\Centering\arraybackslash}O{X}}
\newcolumntype{L}{>{\RaggedRight\arraybackslash}O{X}}
\newcolumntype{R}{>{\RaggedLeft\arraybackslash}O{X}}

% Modulare Tabellenspalten
\newcolumntype{F}{>{\Centering\arraybackslash$}O{X}<{$}} % Zentrierte Formel
\newcolumntype{T}{>{\RaggedRight\arraybackslash}O{X}} % Standard-Text
\newcolumntype{H}{>{\columncolor{\FormulaboxBackColor}\Centering\arraybackslash$}O{X}<{$}} % Formel mit Hintergrund

% Parametrisierte proportionale Spaltentypen
\newcolumntype{f}[1]{>{\hsize=#1\hsize\Centering\arraybackslash$}O{X}<{$}}
\newcolumntype{t}[1]{>{\hsize=#1\hsize\RaggedRight\arraybackslash}O{X}}
\newcolumntype{h}[1]{>{\hsize=#1\hsize\columncolor{\FormulaboxBackColor}\Centering\arraybackslash$}O{X}<{$}}
\newcolumntype{Y}[1]{>{\hsize=#1\hsize\Centering\arraybackslash}O{X}}
\newcolumntype{Z}[1]{>{\hsize=#1\hsize\RaggedRight\arraybackslash}O{X}}
\newcolumntype{Q}[1]{>{\hsize=#1\hsize\RaggedLeft\arraybackslash}O{X}}

\newlength{\ZSFtextMinPad}
\setlength{\ZSFtextMinPad}{1pt}

% Wir binden siunitx ein, es krallt sich aber das S in Tabellen.
% Durch [column=O] verwendet cellspace nun den Bezeichner "O" anstatt "S".
\setlength\cellspacetoplimit{\ZSFtextMinPad}
\setlength\cellspacebottomlimit{\ZSFtextMinPad}
\addparagraphcolumntypes{X,m,p}


% ── Tabellen: Zebra-Ton ──────────────────────────────────────────────
% Default folgt der Kapitelfarbe. Ein fester Ton macht Tabellen über alle
% Kapitel hinweg gleich — sinnvoll, wenn die Kapitelfarbe schon woanders trägt.
% \SetZSFzebraBG{\TableZebraColor}

% ============================================================
% LAYOUT AND SPACING
% ============================================================

\setlength{\parskip}{3.5pt plus 1pt minus 0.5pt}
\setlength{\parindent}{0pt}

\widowpenalty=10000
\clubpenalty=10000
\brokenpenalty=4000
\predisplaypenalty=10000

\raggedcolumns
\newcounter{zsfFirstChap}\setcounter{zsfFirstChap}{1}

% Kompakte Spacing-Skala (4 Stufen)
\newlength{\ZSFspaceXS}          \setlength{\ZSFspaceXS}{1pt}
\newlength{\ZSFspaceS}           \setlength{\ZSFspaceS}{4pt}
\newlength{\ZSFspaceM}           \setlength{\ZSFspaceM}{7pt}
\newlength{\ZSFspaceL}           \setlength{\ZSFspaceL}{10pt}

% Abstand ober-/unterhalb von TextBlocks und Trennlinien
\newlength{\ZSFspaceOuter}        \setlength{\ZSFspaceOuter}{7pt}
\newlength{\ZSFspaceSep}          \setlength{\ZSFspaceSep}{3pt}

\setlength{\ZSFtextMinPad}{\ZSFspaceXS}

\setlength{\columnsep}{\ZSFspaceL} % Konsistent mit Skala
\setlength{\multicolsep}{\ZSFspaceS}
\setlength{\emergencystretch}{2em}
\newlength{\ZSFbreakThreshold}   \setlength{\ZSFbreakThreshold}{3\baselineskip} % Mindesthöhe für BreakIfNotEnoughSpace
\newlength{\ZSFpostChapterNudge} \setlength{\ZSFpostChapterNudge}{0pt}  % Korrektur nach Kapitelbalken (bei Bedarf anpassen)
\newlength{\ZSFtitleBarHeight}    \setlength{\ZSFtitleBarHeight}{9.2pt}  % feste Hoehe aller Titelbalken
\newlength{\ZSFsubsectionBarBeforeSkip}\setlength{\ZSFsubsectionBarBeforeSkip}{\ZSFspaceL}     % Standard-Abstand vor SubsectionBar
\newlength{\ZSFsubsectionAfterChapterGap}\setlength{\ZSFsubsectionAfterChapterGap}{\ZSFspaceXS} % Restabstand wenn SubsectionBar direkt nach ChapterBar steht
\newlength{\ZSFsubsectionAfterGap}       \setlength{\ZSFsubsectionAfterGap}{\ZSFspaceXS}        % Restabstand wenn Box direkt nach SubsectionBar steht
% Bewusste Feinjustierung zwischen XS und S: nicht in Skala, weil der Abstand
% i.d.R. durch \addvspace der SubsectionBar absorbiert wird und nur in den
% Sonderfällen ohne nachfolgende SubsectionBar visuell auftaucht.
\newlength{\ZSFchapterBarAfterSkip}      \setlength{\ZSFchapterBarAfterSkip}{2pt} % Standard-Abstand nach ChapterBar

% Tokens fuer statementbox / procedure (Aussagen & Verfahren)
\newlength{\ZSFstatementBarWidth} \setlength{\ZSFstatementBarWidth}{2pt}   % Breite des linken Akzentstreifens
\newlength{\ZSFstatementLeftPad}  \setlength{\ZSFstatementLeftPad}{\ZSFspaceM} % Abstand Streifen -> Text
\newlength{\ZSFprocStepGap}       \setlength{\ZSFprocStepGap}{\ZSFspaceXS} % Abstand Schritt-Titel -> Beschreibung
\newlength{\ZSFprocItemSep}       \setlength{\ZSFprocItemSep}{\ZSFspaceXS} % Abstand zwischen Schritten

% ============================================================
% TEXT-BOX-BINDUNGEN (Klassen für Abstände)
\newcommand{\ZSFRobustUnskip}{%
  \ifvmode
    \ifdim\lastskip>0pt\relax\vskip-\lastskip\fi
    \ifnum\lastpenalty=0\else\unpenalty\fi
    \ifdim\lastskip>0pt\relax\vskip-\lastskip\fi
  \else
    \ifdim\lastskip>0pt\relax\unskip\fi
    \ifnum\lastpenalty=0\else\unpenalty\fi
    \ifdim\lastskip>0pt\relax\unskip\fi
  \fi
}

% Text, der zur Box DARÜBER gehört (Post-Box Text)
\newcommand{\textNachBox}[1]{%
    \par
    \ZSFRobustUnskip % Löscht zu 100% verlässlich den Skip der Box (ohne manuelles Minus-Rechnen)
    \vspace{\ZSFspaceXS}% Unser exakter Soll-Abstand
    \begingroup
    \setlength{\parskip}{0pt}% Kein zusätzlicher Absatz-Abstand!
    \nopagebreak
    {\small\color{TextNachBoxColor} \noindent #1\par}%
    \endgroup
    \addvspace{\ZSFspaceOuter}% Abstand zum nächsten Element
}

% Text, der zur Box DARUNTER gehört (Pre-Box Text)
\newcommand{\textVorBox}[1]{%
    \par
    \addvspace{\ZSFspaceOuter}%
    \begingroup
    \setlength{\parskip}{0pt}% Kein zusätzlicher Absatz-Abstand!
    {\small\color{TextVorBoxColor} \noindent #1\par}%
    \endgroup
    \nopagebreak
    \vspace{\ZSFspaceXS}% Unser exakter Soll-Abstand
    \vspace{-\ZSFspaceS}% Kompensiert den kommenden Box-Abstand
    \nopagebreak% Kein Spaltenumbruch nach negativem Glue (verhindert Überlappung in multicols)
}

% ============================================================
% TEXT-FORMEL-BINDUNGEN (Innerhalb von formulabox)
% ============================================================

% Titel/Text vor einer Formel (innerhalb der formulabox)
\newcommand{\textVorFormel}[1]{%
    \par
    \begingroup
    \setlength{\parskip}{0pt}% Kein zusätzlicher Absatz-Abstand!
    \noindent\raggedright\textcolor{black!60}{\textbf{#1}}\par
    \endgroup
    \vspace{\ZSFspaceXS}% Unser exakter Soll-Abstand
    \centering
}

% Erklärungstext nach einer Formel (innerhalb der formulabox)
\newcommand{\textNachFormel}[1]{%
    \par
    \vspace{\ZSFspaceXS}% Unser exakter Soll-Abstand
    \begingroup
    \setlength{\parskip}{0pt}% Kein zusätzlicher Absatz-Abstand!
    \noindent\raggedright\textcolor{black!60}{\small #1}\par
    \endgroup
    \centering
}

% ============================================================
% COLORS, CHAPTER PALETTE AND CHAPTER STRUCTURE
% ============================================================

% Palette sortiert nach max. 3-Fenster-Kontrast (CIEDE2000)
% ChapterColor0 bleibt ETH-Rot für den Frontchapter
\definecolor{ChapterColor0}{HTML}{A50D0D}
\definecolor{ChapterColor0Light}{HTML}{EED2D2}
\definecolor{ChapterColor1}{HTML}{5B8C00}
\definecolor{ChapterColor1Light}{HTML}{DDE8B5}
\definecolor{ChapterColor2}{HTML}{7211D4}
\definecolor{ChapterColor2Light}{HTML}{E0D2EE}
\definecolor{ChapterColor3}{HTML}{D44B11}
\definecolor{ChapterColor3Light}{HTML}{EEDAD2}
\definecolor{ChapterColor4}{HTML}{00897B}
\definecolor{ChapterColor4Light}{HTML}{B2DFDB}
\definecolor{ChapterColor5}{HTML}{2B0DA5}
\definecolor{ChapterColor5Light}{HTML}{D8D2EE}
\definecolor{ChapterColor6}{HTML}{A5680D}
\definecolor{ChapterColor6Light}{HTML}{EEE3D2}
\definecolor{ChapterColor7}{HTML}{2E7D32}
\definecolor{ChapterColor7Light}{HTML}{C8E6C9}
\definecolor{ChapterColor8}{HTML}{D411C0}
\definecolor{ChapterColor8Light}{HTML}{EED2EB}
\definecolor{ChapterColor9}{HTML}{1199D4}
\definecolor{ChapterColor9Light}{HTML}{D2E5EE}
\definecolor{ChapterColor10}{HTML}{2BA50D}
\definecolor{ChapterColor10Light}{HTML}{D8EED2}
\definecolor{ChapterColor11}{HTML}{A50D68}
\definecolor{ChapterColor11Light}{HTML}{EED2E3}
\definecolor{ChapterColor12}{HTML}{860DA5}
\definecolor{ChapterColor12Light}{HTML}{E8D2EE}
\definecolor{ChapterColor13}{HTML}{0DA5A5}
\definecolor{ChapterColor13Light}{HTML}{D2EEEE}
\definecolor{ChapterColor14}{HTML}{0D4AA5}
\definecolor{ChapterColor14Light}{HTML}{D2DDEE}
\definecolor{ChapterColor15}{HTML}{D4114B}
\definecolor{ChapterColor15Light}{HTML}{EED2DA}
\definecolor{ChapterColor16}{HTML}{0DA54A}
\definecolor{ChapterColor16Light}{HTML}{D2EEDD}
\definecolor{ChapterColor17}{HTML}{86A50D}
\definecolor{ChapterColor17Light}{HTML}{E8EED2}
\definecolor{ChapterColor18}{HTML}{1124D4}
\definecolor{ChapterColor18Light}{HTML}{D2D5EE}
\definecolor{ChapterColor19}{HTML}{D4C011}
\definecolor{ChapterColor19Light}{HTML}{EEEBD2}
\definecolor{ETHRot}{HTML}{CC0000}

% Semantische Farben für Text-Box-Bindungen
\definecolor{TextVorBoxColor}{HTML}{0A0A1E}   % Sehr dunkles Blau (einleitend)
\definecolor{TextNachBoxColor}{HTML}{0A1A0A}   % Sehr dunkles Grün (ergänzend)

\newcommand{\chaptercolor}{ChapterColor0}
\newcommand{\chaptercolorlight}{ChapterColor0Light}
\newcommand{\ChapterBarColor}{\chaptercolor}
\newcommand{\SubsectionBarColor}{\chaptercolor!75}
\newcommand{\FormulaboxBackColor}{\chaptercolorlight!40}

% Statische Farben für Box-Titel und neutrale Sub-Boxen (unabhängig vom Kapitel)
\definecolor{StaticBoxTitleBack}{HTML}{E2E8F0} % Helles blaugrau (Tailwind Slate-200)
\definecolor{StaticBoxTitleText}{HTML}{0F172A} % Dunkles blaugrau (Tailwind Slate-900)
\definecolor{StaticSubBoxBack}{HTML}{F8FAFC} % Sehr helles grau (Tailwind Slate-50)
\definecolor{StaticSubBoxBorder}{HTML}{CBD5E1} % Mittleres grau (Tailwind Slate-300)

% Dynamische Anpassung an Kapitel, damit Boxen nicht alle gleich aussehen
\newcommand{\BoxTitleColor}{\chaptercolorlight}
\newcommand{\BoxTitleTextColor}{TextVorBoxColor}
\newcommand{\BoxBackColor}{\chaptercolorlight!8}
\newcommand{\BoxFrameColor}{\chaptercolorlight!25}

\newcommand{\GoalConditionBackColor}{StaticSubBoxBack}
\newcommand{\GoalConditionFrameColor}{StaticSubBoxBorder}
\newcommand{\GoalTargetBackColor}{StaticBoxTitleBack}
\newcommand{\GoalTargetFrameColor}{StaticSubBoxBorder}

\newcount\ZSFchapterPaletteIndex
\newcommand{\ZSFchapterPaletteSize}{20}
\newcommand{\ZSFSetChapterPaletteByIndex}[1]{%
  \renewcommand{\chaptercolor}{ChapterColor#1}%
  \renewcommand{\chaptercolorlight}{ChapterColor#1Light}%
}

\newcommand{\SetChapterPaletteByNumber}{%
  \ZSFchapterPaletteIndex=\value{section}%
  \loop
  \ifnum\ZSFchapterPaletteIndex>\numexpr\ZSFchapterPaletteSize-1\relax
    \advance\ZSFchapterPaletteIndex by -\ZSFchapterPaletteSize
  \repeat
  \ifnum\value{section}>\numexpr\ZSFchapterPaletteSize-1\relax
    \PackageWarning{ZSF}{Chapter palette wrapped at section \thesection. Increase palette if unique colors are required}%
  \fi
  \ZSFSetChapterPaletteByIndex{\the\ZSFchapterPaletteIndex}%
}

\newif\ifzsfAfterChapterBar
\zsfAfterChapterBarfalse
\newcommand{\ZSFMarkAfterChapterBar}{\global\zsfAfterChapterBartrue}
% Kollabiert den before-skip der SubsectionBar auf \ZSFsubsectionAfterChapterGap,
% wenn sie direkt nach einer ChapterBar steht. Reset nach Verwendung.
\newcommand{\ZSFCollapseAfterChapterBar}{%
  \ifzsfAfterChapterBar
    \vspace*{\dimexpr\ZSFsubsectionAfterChapterGap-\ZSFsubsectionBarBeforeSkip\relax}%
    \global\zsfAfterChapterBarfalse
  \fi
}

\makeatletter
\newcommand{\ZSFPostChapterSpacing}{\vspace*{\ZSFpostChapterNudge}}
\newcommand{\StartChapter}[2][]{%
  \ifnum\value{zsfFirstChap}=1\setcounter{zsfFirstChap}{0}\fi
  \refstepcounter{section}
  \SetChapterPaletteByNumber
  \setcounter{subsection}{0}
  \if\relax\detokenize{#1}\relax\else\label{#1}\fi
  \ChapterBar{\thesection. #2}
  \ZSFMarkAfterChapterBar
  \ZSFPostChapterSpacing
}
\makeatother

% Variante: Kapitel auf neuer Spalte beginnen (ersetzt rohes \columnbreak)
\newcommand{\StartChapterOnNewColumn}[2][]{%
  \columnbreak
  \StartChapter[#1]{#2}%
}

\newcommand{\StartFrontChapter}[1]{
  \ifnum\value{zsfFirstChap}=1\setcounter{zsfFirstChap}{0}\else\columnbreak\fi
  \renewcommand{\chaptercolor}{ChapterColor0}
  \renewcommand{\chaptercolorlight}{ChapterColor0Light}
  \setcounter{subsection}{0}
  \ChapterBar{#1}
  \ZSFMarkAfterChapterBar
  \ZSFPostChapterSpacing
}

% ============================================================
% 50_typography_semantics.tex — Schriftmakros für Box-Typen & semantische Inhaltsmakros
% ============================================================

% ── Die eine Rechenstelle für jede benannte Schriftrolle ─────────────
% Jede Grösse unten wird GERECHNET, nicht geschrieben: Sie ist ihr bp-Mass
% mal dem Rollenfaktor mal der Grundgrösse der ZSF (\ZSFBodyScale aus
% 30_layout_spacing, 1 solange \ZSFBodyFontSize nicht steht). Über \the und
% damit in pt, weil \fontsize sein Argument sonst als Zahl ohne Einheit lesen
% müsste — der Faktor kommt so als fertiges Mass an und nicht als Term.
%
% \ZSFBodyScale steht hier fest verdrahtet und nicht als vierter Parameter:
% Es gibt keine Schriftrolle im Textkörper, die der Grundgrösse NICHT folgen
% soll. Wäre sie ein Argument, wäre die Frage bei jeder neuen Rolle wieder
% offen — und genau eine Rolle bekäme sie irgendwann anders beantwortet.
\makeatletter
\newlength{\ZSF@fontSize}
\newlength{\ZSF@fontSkip}
\newcommand{\ZSF@scaledFont}[3]{%
  \setlength{\ZSF@fontSize}{\ZSFBodyScale\dimexpr #1\dimexpr #2\relax\relax}%
  \setlength{\ZSF@fontSkip}{\ZSFBodyScale\dimexpr #1\dimexpr #3\relax\relax}%
  \fontsize{\the\ZSF@fontSize}{\the\ZSF@fontSkip}\selectfont}

% Explizite bp-Grössen (wie Analysis) für präzise, font-unabhängige Kontrolle.
% Der Rollenfaktor ist 1, wo die Rolle keinen eigenen Regler hat — nur die
% Diagramm-Beschriftung bringt mit \ZSFDiagramLabelScale einen mit.
\newcommand{\ZSFfontChapter}{\sffamily\bfseries\ZSF@scaledFont{1}{8.1bp}{9.1bp}}
\newcommand{\ZSFfontSection}{\sffamily\bfseries\ZSF@scaledFont{1}{8.1bp}{9.1bp}}
\newcommand{\ZSFfontSubsubsection}{\sffamily\bfseries\ZSF@scaledFont{1}{7.2bp}{8.0bp}}
\newcommand{\ZSFfontBoxTitle}{\sffamily\bfseries\ZSF@scaledFont{1}{8.1bp}{9.1bp}}
% chktex 6
\newcommand{\ZSFfontNote}{\sffamily\itshape\ZSF@scaledFont{1}{6.7bp}{7.5bp}}
% Rechtsbuendiges Meta-/Kategorie-Tag im Titelbalken (vgl. \ZSFTitleTag in
% styles/60_boxes.tex). Etwas kleiner und kursiv, damit es den Titel begleitet
% statt mit ihm zu konkurrieren.
\newcommand{\ZSFfontTitleTag}{\sffamily\itshape\ZSF@scaledFont{1}{6.4bp}{7.0bp}}
\newcommand{\ZSFfontTitleHeader}{\sffamily\bfseries\ZSF@scaledFont{1}{11bp}{12.5bp}}
\makeatother

% ── Seitenmöbel: bewusst OHNE \ZSFBodyScale ──────────────────────────
% Fusszeile und Identity-Marker gehören zur SEITE und nicht zur Textspalte.
% Die Seitenzahl am Fuss ist eine feste physische Grösse — sie schrumpft nicht
% mit, nur weil der Text der ZSF enger gesetzt ist, und der 4pt-Marker wäre
% eine Stufe kleiner nicht mehr lesbar. Deshalb rohes \fontsize: Das ist hier
% die Aussage »absolut gemeint«, nicht ein vergessener Regler.
\newcommand{\ZSFfontFooter}{\sffamily\bfseries\fontsize{9.5pt}{9.5pt}\selectfont}
\newcommand{\ZSFfontIdentityMarker}{\fontsize{4pt}{4pt}\selectfont}

% Schrift/Farbe für Code-Kommentare — nur für das optionale Code-Modul
% (styles/67_code_comments.tex) benötigt. ce_green stammt aus
% styles/40_colors_structure.tex. Schadlos wenn das Modul deaktiviert ist.
\newcommand{\ZSFfontCodeComment}{\ttfamily\color{ce_green}}

% Semantische Inhalts-Fontgrössen — die zwei Stufen des Reglers `font`
% (`normal` | `dense`). Kapitel wählen sie nicht direkt, sondern über den
% Regler an der Box (styles/60_boxes.tex) oder an der ZSFtable
% (styles/20_tables.tex). `dense` ist die eine Stufe kleiner: lange Register,
% Vergleichsmatrizen, gedrängte Aufzählungen.
%
% EIN Paar für beide Bausteine, nicht je eines: »eine Stufe kleiner« ist
% dieselbe Entscheidung, ob sie eine Tabelle oder einen Boxinhalt trifft. Mit
% zwei Paaren müsste jede Änderung am einen im anderen nachgezogen werden
% (styles/README → Vereinigungstest).
%
% Beide holen die Display-Skips zurück: \normalsize und \footnotesize setzen
% sie auf die Klassenwerte, und die ZSF-Entscheidung dazu ist eine andere
% (\ZSFApplyDisplaySkips in 30_layout_spacing). Ohne diesen Nachsatz wären
% Formeln in jeder Box luftiger als im Fliesstext — lautlos, weil beides
% für sich korrekt aussieht.
% ── Der Rollenfaktor des Bausteininhalts ────────────────────────────
% Das Gegenstück zu \ZSFFlowTextScale, und aus demselben Grund: Wie laut steht
% der Inhalt der Bausteine — Boxtext, Tabellenzelle, Formel — neben den
% Titelbalken, die ihn gliedern?
%
% Ohne diesen Faktor gab es dafür nur \ZSFBodyFontSize, und das nimmt die
% Balken MIT. Der Zustand »Balken auf voller Grösse, Inhalt eine Stufe
% darunter« war damit nicht ausdrückbar — obwohl er der natürliche ist: Die
% Balkengrössen stehen absolut in bp (oben) und sind auf die Spaltenbreite
% abgestimmt, der Inhalt dagegen muss in eine Tabellenzelle passen. Das sind
% zwei Fragen, und sie hatten einen gemeinsamen Regler.
%
% Gemessen, woran das auffiel: Eine Formelzelle einer zweispaltigen Tabelle
% ist 86.2 pt breit. \sin(\arctan(x)) = x/\sqrt{1+x^2} braucht bei einem
% Inhalt von 8 pt 86.3 pt und bricht um; bei 7 pt sind es 75.9 pt. Wer das
% ueber \ZSFBodyFontSize loest, schrumpft die Balken von 8.1 auf 7.1 bp mit
% und aendert das Satzbild an einer Stelle, an der nichts falsch war.
%
% Angewandt NACH der Klassenstufe, nicht statt ihr: \normalsize und
% \footnotesize tragen \ZSFBodyScale bereits (30_layout_spacing rechnet die
% Stufen um, sobald \ZSFBodyFontSize steht). Der Rollenfaktor multipliziert
% nur noch — sonst zaehlte die Grundgroesse doppelt.
%
% Ein Faktor fuer beide Stufen, nicht zwei: »eine Stufe kleiner« (`font=dense`)
% ist eine Entscheidung PRO STELLE und bleibt es; wie laut der Inhalt insgesamt
% steht, ist eine Entscheidung pro ZSF. Zwei Faktoren wuerden die beiden Fragen
% vermischen (styles/README → Vereinigungstest).
\providecommand{\ZSFContentScale}{1}
\makeatletter
\newcommand{\ZSF@stepScale}[1]{%
  \setlength{\ZSF@fontSize}{#1\dimexpr\f@size pt\relax}%
  \setlength{\ZSF@fontSkip}{#1\dimexpr\f@baselineskip\relax}%
  \fontsize{\the\ZSF@fontSize}{\the\ZSF@fontSkip}\selectfont}
% Bei Faktor 1 gar nicht erst umschalten: Ein \fontsize auf denselben Wert
% kostet eine Fontauswahl und kann durch die pt-Rundung minimal danebenliegen.
\newcommand{\ZSF@applyContentScale}{%
  \ifdim\ZSFContentScale pt=1pt\else\ZSF@stepScale{\ZSFContentScale}\fi}
\newcommand{\ZSFfontContent}{\normalsize\ZSF@applyContentScale\ZSFApplyDisplaySkips}
\newcommand{\ZSFfontContentDense}{\footnotesize\ZSF@applyContentScale\ZSFApplyDisplaySkips}
\makeatother

% ── Freier Fliesstext: die dritte Inhaltsrolle ───────────────────────
% `runintext` ist der einzige Inhalts-Baustein OHNE eigene Schriftrolle — er
% erbt schlicht die Umgebungsgrösse. Damit war der Fliesstext das einzige
% Stück Satz, das man nur über einen rohen \fontsize erreichen konnte, und
% genau dort landete historisch der Notausgang in main.tex.
%
% Als eigene Rolle ist es eine Entscheidung statt eines Nebeneffekts: Wie laut
% steht die verbindende Prosa neben dem Bausteininhalt, der sie umgibt? In
% einer box-lastigen ZSF trägt die Box den prüfungsrelevanten Inhalt und der
% Fliesstext verbindet — dann darf er eine Spur leiser sein. In einer
% textlastigen ist es umgekehrt und der Faktor bleibt auf 1.
%
% Gerechnet wird von der KLASSENGRÖSSE aus, nicht vom aktuellen Font: So
% multiplizieren sich Grundgrösse und Rollenfaktor sauber (\ZSF@scaledFont
% legt \ZSFBodyScale selbst davor), statt sich zu verdoppeln, wenn die Rolle
% in einem bereits umgeschalteten Kontext aufgerufen wird.
\providecommand{\ZSFFlowTextScale}{1}
\makeatletter
\newcommand{\ZSFfontFlow}{%
  \ZSF@scaledFont{\ZSFFlowTextScale}{\ZSF@classFontSize pt}{\ZSF@classBaselineSkip}%
  \ZSFApplyDisplaySkips}
\makeatother

%
%
%
%
%
%
%
%
%
\newcommand{\ZSFfontTableHead}{\bfseries}
% Marke einer ZSFlist und Beschriftung eines \ZSFsep-Blockwechsels — beide
% sind dieselbe Rolle: eine kurze fette Beschriftung innerhalb eines Bausteins.
\newcommand{\ZSFfontBlockLabel}{\bfseries}
% Text, der per \textVorBox/\textNachBox an eine Box gebunden wird.
\newcommand{\ZSFfontBoxBinding}{\small}
% Inline-Pille für Stolperfallen (\ZSFdanger).
\newcommand{\ZSFfontDangerTag}{\sffamily\bfseries}
%
%
%
\newcommand{\ZSFfontRefLocator}{\upshape\bfseries}

%
%
%
%
%
\newcommand{\ZSFBoxNote}[1]{{\ZSFfontNote\color{MutedTextColor}#1}}

%
%
%
%
%
%
%
%
%
%
%
%
%
%
%
%
%
%
%
%
%
%
%
%
%
%
\providecommand{\ZSFDiagramLabelScale}{1}
% Gerechnet wird oben in \ZSF@scaledFont — dieselbe Stelle wie für jede andere
% Schriftrolle. Der Rollenfaktor \ZSFDiagramLabelScale multipliziert dort die
% Grundgrösse der ZSF: Ein diagrammlastiges Blatt dreht seine Beschriftungen
% kleiner, ohne dass sie aufhören, dem Satz der ZSF zu folgen.
\makeatletter
\newcommand{\ZSFfontDiagramLabel}{%
  \sffamily\ZSF@scaledFont{\ZSFDiagramLabelScale}{5.8bp}{6.4bp}}
\newcommand{\ZSFfontDiagramLabelSmall}{%
  \sffamily\ZSF@scaledFont{\ZSFDiagramLabelScale}{4.6bp}{5.1bp}}
\makeatother

%
%
%
%
%
%
%
%
%
%
%
%
%
%
%
%
%
%
%
%
%
%
%
%
\NewDocumentCommand{\ZSFkeyword}{s o m}{\ZSFkeywordStyle{#3}}

%
%
%
%
%
\newcommand{\ZSFlabel}[1]{\ZSFlabelStyle{#1}}

% Die beiden Darstellungs-Hooks der Marker — getrennt, weil \ZSFkeyword und
% \ZSFlabel verschiedene Dinge bedeuten und eine ZSF den Unterschied ändern
% können soll, ohne ein einziges Kapitel anzufassen.
%
% BEIDE schlicht fett, ohne Farbe. \ZSFkeyword stand bis hierher fett in der
% KAPITELFARBE, mit dem Argument, der primäre Scan-Anker müsse sich von jeder
% anderen Fettung abheben. Das Argument ist am Satz gescheitert: In laufendem
% Text bringt eine eingefärbte Vokabel Unruhe in den Absatz, ohne dass die
% Farbe etwas sagt — sie wiederholt nur, in welchem Kapitel man ohnehin gerade
% ist. Wer den Absatz liest, wird gestört; wer ihn überfliegt, findet die
% Fettung auch ohne Farbe.
%
% Farbe bleibt damit im System das, was sie vorher schon war und wofür sie
% trägt: Kapitel-Identität auf Flächen (Balken, Akzente), Wegweiser zum ZIEL
% eines Verweises (\ZSFref) und semantische Zuordnung in Formeln (\ZSFmhl*).
% Als blosse Betonung im Fliesstext ist sie ausdrücklich nicht vorgesehen.
%
% Die Unterscheidung der beiden Marker liegt jetzt allein in der Bedeutung —
% \ZSFkeyword ist ein Fachbegriff und landet im Register, \ZSFlabel ist eine
% neutrale Beschriftung und tut das nicht. Das ist der Unterschied, der beim
% Schreiben zählt; er war nie der Farbe zu entnehmen.
\newcommand{\ZSFkeywordStyle}[1]{\textbf{#1}}
\newcommand{\ZSFlabelStyle}[1]{\textbf{#1}}

%
%
%
%
%
%
%
\newcommand{\ZSFScriptRef}[1]{\hfill\ZSFInk{\ScriptRefColor}{\ZSFfontBoxBinding (S.#1)}}

% \ZSFhl{...} — Hervorhebung des EINEN prüfungskritischsten Fakts pro Box
% (fett + unterstrichen). Bewusst sparsam: hoechstens einmal pro Block.
% Klare Rollentrennung:
%   * Fachbegriff/Scan-Anker        -> \ZSFkeyword
%   * Stolperfalle/kritische Ausnahme -> \ZSFdanger
%   * Folgerung                      -> \ZSFconclusion
%   * sonst die EINE Schlüsselaussage -> \ZSFhl
% \uline statt \underline: \underline steckt seinen Inhalt in eine
% unzerbrechliche Box — jede Hervorhebung, die länger als eine Zeile ist,
% sprengt damit die 50-mm-Spalte. \ZSFhl markiert aber laut Zweck eine ganze
% Aussage, nicht ein Wort. \uline bricht mit dem Absatz um.
% Fett als Deklaration, nicht als Argument von \uline: ein \textbf{...} als
% Argument macht den Inhalt für ulem wieder zu einer Einheit.
%
% Die Darstellung liegt wie bei \ZSFkeyword und \ZSFlabel in einem eigenen
% Hook. Ohne ihn erwischte »die Hervorhebung dieser ZSF umstellen« zwei von
% drei Markern — der dritte war der einzige mit fest verdrahteter Erscheinung,
% und man sah es erst am Satz (styles/README → Auszeichnungs-Hooks).
\newcommand{\ZSFhlStyle}[1]{{\bfseries\uline{#1}}}
\newcommand{\ZSFhl}[1]{\ZSFhlStyle{#1}}

% Papierorientierte Referenzrollen:
%   * \ZSFref{label}: hilfreicher Querverweis mit farbiger Abschnittsnummer.
%   * \ZSFsectionref{label}: kompakte Zielnummer für lokale Router/Tabellen.
% Ziel-Label via optionales Argument von \SubsubsectionBar/\SubsectionBar/\StartChapter setzen:
%   \SubsectionBar[sec:matrix-inverse]{Inverse}
% Farbe = Palette-Farbe des ZIELkapitels (Wegweiser-Prinzip): wird automatisch
% aus der Kapitelnummer der Referenz abgeleitet ("6.6" -> ChapterColor6).
% Bei Kapiteln jenseits der Palette wird derselbe Modulo-Lookup wie bei
% \StartChapter genutzt, damit gewrappte Kapitel und Verweise konsistent sind.
% Aufrecht + fett statt kursiv für maximale Lesbarkeit bei 6.7bp.
\makeatletter
\def\ZSF@refExtractFirst#1#2\@nil{#1}
\def\ZSF@refExtractChap#1.#2\@nil{#1}
% Bestimmt \ZSF@refTargetColor aus dem Label; Fallback auf die aktive
% Kapitelfarbe im ersten latexmk-Pass (Label noch undefiniert).
\newcommand{\ZSF@refComputeTargetColor}[1]{%
  \@ifundefined{r@#1}{%
    \edef\ZSF@refTargetColor{\chaptercolor}%
  }{%
    \edef\ZSF@refTargetNum{\expandafter\expandafter\expandafter
      \ZSF@refExtractFirst\csname r@#1\endcsname\@nil}%
    \edef\ZSF@refTargetChap{\expandafter\ZSF@refExtractChap\ZSF@refTargetNum.\@nil}%
    \ZSFSetPaletteColorNameFromNumber{\ZSF@refTargetChap}{\ZSF@refTargetColor}%
  }%
}
% Die Zielfarbe läuft über \ZSFInk. Ohne das steht ein Verweis auf das eigene
% Kapitel im Titel einer \SubsectionBar als ChapterColorN auf \chaptercolor!75 —
% derselbe Farbton bei voller Sättigung, also unlesbar, ohne dass Build oder
% Linter etwas melden. Auf einer Fläche mit eigener Kontrastfarbe erbt der
% Verweis diese und bleibt lesbar; der Farb-Wegweiser gilt dort, wo er auch
% ankommt — im Boxinhalt (40_colors_structure → Ink-Vertrag).
\newcommand{\ZSFref}[1]{%
  \begingroup
  \ZSF@refComputeTargetColor{#1}%
  \hyperref[#1]{{\ZSFfontNote\ZSFfontRefLocator
    \ZSFInk{\ZSF@refTargetColor}{($\rightarrow$\,\ref*{#1})}}}%
  \endgroup
}
\newcommand{\ZSFsectionref}[1]{%
  \begingroup
  \ZSF@refComputeTargetColor{#1}%
  \hyperref[#1]{{\ZSFfontRefLocator\ZSFInk{\ZSF@refTargetColor}{\ref*{#1}}}}%
  \endgroup
}
\makeatother

\newcommand{\ZSFconclusion}[1]{%
  \ZSFInterlude{\ZSFspaceS}%
  \noindent$\Rightarrow$ #1%
  \ZSFInterlude{\ZSFspaceS}%
}



% ── Marker-Darstellung ───────────────────────────────────────────────
% Default: \ZSFkeyword und \ZSFlabel beide schlicht fett, ohne Farbe.
% Inline-Betonung im Fliesstext trägt bewusst KEINE Kapitelfarbe — sie brächte
% Unruhe in den Absatz und sagte nichts, was die Seite nicht ohnehin zeigt.
% Farbe bleibt den Flächen, den Verweisen und den Formel-Markern vorbehalten.
% Umstellen heisst GENAU diese eine Zeile, nie ein Kapitel:
% \renewcommand{\ZSFkeywordStyle}[1]{\textbf{#1}}
% \renewcommand{\ZSFlabelStyle}[1]{\textbf{#1}}
% Dritter Marker derselben Familie: die Hervorhebung der Kernaussage.
% \renewcommand{\ZSFhlStyle}[1]{{\bfseries #1}}

% ── Grössenfarben: Farbe als Identität einer Grösse ──────────────────
% Für Fächer, in denen dieselben Grössen über das ganze Dokument hinweg in
% Formeln auftauchen (Kraft, Moment, Spannung, Durchmesser …).
%
% In den Theoriekapiteln arbeitet Analysis 2 positional mit \ZSFmhlA..D
% innerhalb einer Herleitung — f, x, t sind Variablen und keine Grössenfamilien.
%
% Die Integraltabelle im Anhang ist der eine Fall, wo Identität gilt: Über
% mehrere Dutzend Stammfunktionen hinweg bezeichnen `a` und `b` durchgehend
% dieselben zwei freien Parameter. Wer eine Formel überträgt, muss auf einen
% Blick sehen, welcher Parameter wo landet — genau die Zusage dieses Systems.
% Neutrale Namen mit Absicht: In \int \frac{1}{(x-a)(x-b)} sind beide
% gleichrangige Nullstellen, eine Benennung wie »Streckung/Verschiebung«
% wäre dort schlicht falsch.
\ZSFDeclareQuantity{ParameterA}{9}   % danach: \ZSFqParameterA{a}  (Himmelblau)
\ZSFDeclareQuantity{ParameterB}{6}   % danach: \ZSFqParameterB{b}  (Grün)

% ── Eigener Ton: eine wiederkehrende Farbwelt fürs Fach ──────────────
% `tone` hat drei mitgelieferte Werte (chapter, neutral, warn). Braucht das
% Fach eine WIEDERKEHRENDE visuelle Kategorie, die keiner davon trägt (ein
% Beispiel-Block, ein Merksatz), ist die Antwort ein weiterer Ton — keine
% weitere Box. Nach der Deklaration gilt `tone=beispiel` auf JEDER Box.
% \ZSFDeclareTone{beispiel}{ ... }

\input{styles/55_readability}

% ── Fliesstext: Flattersatz ──────────────────────────────────────────
% \ZSFReadableBodyOn ist der Default, weil die ~50 mm schmalen Spalten im
% Blocksatz entweder trennen oder Wortzwischenräume aufblähen müssen.
% \ZSFReadableBodyOff             % zurück zu Blocksatz

% ============================================================
% BOXES, CHAPTER/SUBSECTION BARS, AND HELPERS
% ============================================================

\newcommand{\BreakIfNotEnoughSpace}{%
  \par\penalty0\vspace{0pt plus \ZSFbreakThreshold}\penalty-100\vspace{0pt plus -\ZSFbreakThreshold}\relax
}

% Einheitlicher Top-Spacing-Kontrakt fuer Titelbalken.
% Reihenfolge ist kritisch: \addvspace MUSS vor \BreakIfNotEnoughSpace stehen,
% damit es per Max-Semantik mit dem after-skip der Vorbox kollabiert. Stuende
% es danach, faende \addvspace nur noch die Stretch-Reste (lastskip natural=0)
% und wuerde additiv arbeiten -> Vorbox-After-Skip + Bar-Before-Skip stapeln
% sich, was zu sichtbar variablem Abstand ueber Titelbalken fuehrt.
% Tcolorbox-Definitionen, die diesen Helper aufrufen, MUESSEN before skip=0pt
% setzen, damit das Spacing nicht doppelt eingerechnet wird.
\newcommand{\ZSFTitleBarTopSpacing}[1]{%
  \par
  \addvspace{#1}%
  \BreakIfNotEnoughSpace
}

\newtcolorbox{chapterbar}[1][] {
  enhanced jigsaw,
  breakable=false,
  colback=\ChapterBarColor,
  colframe=\ChapterBarColor,
  boxrule=0pt,
  boxsep=0pt,
  left=\ZSFspaceS,
  right=\ZSFspaceS,
  top=\ZSFspaceS,
  bottom=\ZSFspaceS,
  before skip=\ZSFspaceL,
  after skip=\ZSFchapterBarAfterSkip,
  fontupper=\ZSFfontChapter,
  #1
}
\newcommand{\ChapterBar}[1]{
  \begin{chapterbar}
  \textcolor{white}{#1}
  \end{chapterbar}
}

% ------------------------------------------------------------
% Dokument-Titel-Masthead (einmalig, vor dem ersten Kapitel)
% Kompakte Variante fuer Analysis: knapper Platz im 4-Spalten-Layout.
% ------------------------------------------------------------
\newtcolorbox{zsftitleheader}[1][]{
  enhanced jigsaw,
  breakable=false,
  colback=ChapterColor0,
  colframe=ChapterColor0,
  boxrule=0pt,
  boxsep=0pt,
  left=\ZSFspaceM,
  right=\ZSFspaceM,
  top=\ZSFspaceS,
  bottom=\ZSFspaceS,
  before skip=0pt,
  after skip=\ZSFspaceXS,
  halign=center,
  #1
}
\newcommand{\ZSFTitleHeader}[2]{%
  \penalty-500
  \begin{zsftitleheader}
    \sffamily
    {\color{white}\fontsize{11bp}{12.5bp}\bfseries\selectfont #1\par}%
    \vspace{1pt}%
    {\color{white!88}\fontsize{6.7bp}{7.5bp}\itshape\selectfont von~#2\par}%
  \end{zsftitleheader}
  \penalty10000
}

\newtcolorbox{subsectionbar}[1][] {
  enhanced jigsaw,
  breakable=false,
  colback=\SubsectionBarColor,
  colframe=\SubsectionBarColor,
  boxrule=0pt,
  boxsep=0pt,
  left=\ZSFspaceS,
  right=\ZSFspaceS,
  top=\ZSFspaceS,
  bottom=\ZSFspaceS,
  % before skip wird vom Wrapper (\ZSFTitleBarTopSpacing) verwaltet.
  before skip=0pt,
  after skip=\ZSFsubsectionAfterGap,
  fontupper=\ZSFfontSection,
  #1
}
\makeatletter
\newcommand{\SubsectionBar}{\@ifstar{\SubsectionBarStar}{\SubsectionBarNumbered}}
\newcommand{\SubsectionBarNumbered}[2][]{
  \ZSFTitleBarTopSpacing{\ZSFsubsectionBarBeforeSkip}
  \penalty-500
  \ZSFCollapseAfterChapterBar
  \refstepcounter{subsection}
  \if\relax\detokenize{#1}\relax\else\label{#1}\fi
  \begin{subsectionbar}
  \textcolor{white}{\thesubsection. #2}
  \end{subsectionbar}
  \penalty10000
  \ZSFMarkAfterSubsectionBar
}
\newcommand{\SubsectionBarStar}[2][]{
  \ZSFTitleBarTopSpacing{\ZSFsubsectionBarBeforeSkip}
  \penalty-500
  \ZSFCollapseAfterChapterBar
  \if\relax\detokenize{#1}\relax\else\label{#1}\fi
  \begin{subsectionbar}
  \textcolor{white}{#2}
  \end{subsectionbar}
  \penalty10000
  \ZSFMarkAfterSubsectionBar
}
% Variante: Abschnitt auf neuer Spalte beginnen (ersetzt rohes \columnbreak)
\newcommand{\SubsectionBarOnNewColumn}{\@ifstar{\SubsectionBarOnNewColumnStar}{\SubsectionBarOnNewColumnNumbered}}
\newcommand{\SubsectionBarOnNewColumnNumbered}[2][]{%
  \columnbreak
  \SubsectionBarNumbered[#1]{#2}%
}
\newcommand{\SubsectionBarOnNewColumnStar}[2][]{%
  \columnbreak
  \SubsectionBarStar[#1]{#2}%
}
\makeatother

% Gemeinsamer Basis-Stil fuer alle Titelleisten
\tcbset{zsftitlebar/.style={
  fonttitle=\ZSFfontBoxTitle\strut,
  halign title=center,
  title style={minimum height=\ZSFtitleBarHeight},
  lefttitle=0pt,
  righttitle=0pt,
  toptitle=\ZSFspaceXS,
  bottomtitle=\ZSFspaceXS,
  before title={},
  after title={},
}}

% Rechtsbuendiges Meta-/Kategorie-Tag im Titelbalken einer Titelbox.
% Im Titel-Argument verwenden, z.B.:
%   \begin{defbox}[Satz von Fubini\ZSFTitleTag{Integration}] ... \end{defbox}
\newcommand{\ZSFTitleTag}[1]{\hfill{\ZSFfontTitleTag #1}}

% Gemeinsamer Basis-Stil fuer alle Panel-Boxen mit Titelbalken
\tcbset{zsfpanelbox/.style={
  enhanced jigsaw,
  colback=white,
  colbacktitle=\BoxTitleColor,
  colframe=\BoxTitleColor,
  coltitle=\BoxTitleTextColor,
  boxrule=0.5pt,
  titlerule=0pt,
  boxsep=0pt,
  arc=1pt, outer arc=1pt,
  before skip=\ZSFBoxBeforeSkip,
  after skip=\ZSFspaceS,
  before={\ZSFConsumeAfterSubsectionBar},
  zsftitlebar,
}}

% Gemeinsamer Basis-Stil fuer alle Titelboxen (defbox, tablebox, figbox)
\tcbset{zsftitlebox/.style={
  enhanced jigsaw,
  breakable=false,
  colback=\BoxBackColor,
  colbacktitle=\BoxTitleColor,
  colframe=\BoxFrameColor,
  coltitle=\BoxTitleTextColor,
  boxrule=0.5pt,
  titlerule=0pt,
  boxsep=0pt,
  before skip=\ZSFBoxBeforeSkip,
  after skip=\ZSFspaceS,
  before={\ZSFConsumeAfterSubsectionBar},
  zsftitlebar,
  before upper={\parskip=\ZSFspaceS\relax},
}}

\newtcolorbox{defbox}[1][]{zsftitlebox, before skip=\ZSFBoxBeforeSkip, left=\ZSFspaceS, right=\ZSFspaceS, top=\ZSFspaceS, bottom=\ZSFspaceS, title={#1}}
\newtcolorbox{tablebox}[1][]{zsftitlebox, unbreakable, title={#1}}
\newtcolorbox{figbox}[1][]{zsftitlebox, unbreakable, title={#1}}

% Randlose, modulare Tabelle (für zentrierte Formeln, nützt volle Box aus)
\newtcolorbox{fulltablebox}[2][]{
  zsfpanelbox,
  colframe=black!20, % Hellerer Rahmen für die äußere Tabellenbox
  breakable=false, % Tabellen lieber nicht umbrechen
  title={#1},
  boxsep=0pt, left=0pt, right=0pt, top=\ZSFtextMinPad, bottom=\ZSFtextMinPad,
  tabularx*={\arrayrulecolor{\ZSFtableRuleColor}\setlength{\arrayrulewidth}{\ZSFtableRuleWidth}}{#2},
  zsftitlebar,      % Restore title spacing broken by tabularx defaults
  boxrule=0.5pt     % Restore boxrule broken by tabularx defaults
}

% Falls kein Titel, aber volle Breite & zentriert gesucht wird (Randlose Box)
\newtcolorbox{fulllessbox}[1]{
  enhanced jigsaw,
  breakable=false,
  colback=white,
  colframe=white, % Unsichtbarer Rand
  boxrule=0pt,
  arc=0pt, outer arc=0pt,
  before skip=\ZSFspaceS,
  after skip=\ZSFspaceS,
  boxsep=0pt, left=0pt, right=0pt, top=\ZSFtextMinPad, bottom=\ZSFtextMinPad,
  tabularx*={\arrayrulecolor{\ZSFtableRuleColor}\setlength{\arrayrulewidth}{\ZSFtableRuleWidth}}{#1},
  boxrule=0pt       % Restore boxrule broken by tabularx defaults
}

\newenvironment{longformula}{\[\begin{aligned}}{\end{aligned}\]}
\newtcolorbox{formulabox}[1][]{
  enhanced jigsaw,
  colback=\FormulaboxBackColor,
  colframe=black,  % 1px Rand schwarz
  boxrule=0.5pt,   % Ca. 1px Strichdicke in LaTeX
  left=\ZSFspaceS,
  right=\ZSFspaceS,
  top=\ZSFspaceS,
  bottom=\ZSFspaceS,
  before skip=\ZSFBoxBeforeSkip,
  after skip=\ZSFspaceS,
  before={\ZSFConsumeAfterSubsectionBar},
  before upper={\parskip=\ZSFspaceS\centering},
  after upper={\par},
  middle=\ZSFspaceSep,
  #1
}

\newtcolorbox{goalbox}[1][]{
  zsfpanelbox,
  breakable=false,
  colback=white,
  left=\ZSFspaceS,
  right=\ZSFspaceS,
  top=\ZSFtextMinPad,
  bottom=\ZSFtextMinPad,
  title={#1},
  before upper={\parskip=\ZSFspaceS\centering},
  after upper={\par},
}

\newtcbox{\GoalConditionBox}{%
  on line,
  enhanced,
  boxrule=1.5pt,
  colback=white,
  colframe=\chaptercolorlight!22,
  borderline={0.2pt}{0pt}{black},
  arc=1pt,
  left=\ZSFspaceXS,
  right=\ZSFspaceXS,
  top=\ZSFspaceXS,
  bottom=\ZSFspaceXS,
  nobeforeafter,
  tcbox raise base,
}

\newtcbox{\GoalTargetBox}{%
  on line,
  enhanced,
  boxrule=0.5pt,
  colback=\chaptercolorlight!22,
  colframe=\chaptercolor,
  arc=1pt,
  left=\ZSFspaceXS,
  right=\ZSFspaceXS,
  top=\ZSFspaceXS,
  bottom=\ZSFspaceXS,
  nobeforeafter,
  tcbox raise base,
}

\newcommand{\GoalCondition}[1]{\par\centering\GoalConditionBox{\ZSFfontNote #1}\par}
\newcommand{\GoalStep}[1]{\par\centering$\displaystyle #1$\par}
\newcommand{\GoalTarget}[1]{\par\centering\GoalTargetBox{$\displaystyle #1$}\par}
\newcommand{\ZSFDerivationCase}[3]{%
  \GoalCondition{#1}%
  \GoalStep{#2 =}%
  \GoalTarget{#3}%
}

\newcommand{\ZSFDerivationInlineCase}[3]{%
  \GoalCondition{#1}%
  \GoalStep{#2 = \GoalTargetBox{$\displaystyle #3$}}%
}

% Kompaktere Variante mit reduziertem vertikalen Abstand zwischen Bedingung und Umformung
\newcommand{\ZSFDerivationInlineCaseCompact}[3]{%
  {%
    \setlength{\parskip}{0pt}%
    \par\raggedright%
    \tcbox[
      on line,
      enhanced,
      boxrule=1.5pt,
      colback=white,
      colframe=\chaptercolorlight!22,
      borderline={0.2pt}{0pt}{black},
      arc=1pt,
      boxsep=0pt,
      left=0pt,
      right=0pt,
      top=0pt,
      bottom=0pt,
      nobeforeafter,
      tcbox raise base
    ]{\ZSFfontNote #1}\par
    \vspace{0pt}%
    \centering$\displaystyle #2 = \GoalTargetBox{$\displaystyle #3$}$\par
    \vspace{-0.5pt}%
  }%
}

\newcommand{\GoalCard}[4][]{%
  \begin{goalbox}[#1]
    \GoalCondition{#2}
    #3
    \GoalTarget{#4}
  \end{goalbox}%
}

\newcommand{\GoalPair}[8]{%
  \begin{tcbraster}[raster columns=2, raster column skip=\ZSFspaceM, raster row skip=0pt, raster equal height]
    \GoalCard[#1]{#2}{#3}{#4}
    \GoalCard[#5]{#6}{#7}{#8}
  \end{tcbraster}%
}

% \formulanote: Note-Text unter Formeln - mit parskip-Kompensation
\newcommand{\formulanote}[1]{%
  \par\vspace{-\parskip}\vspace{\ZSFspaceS}%
  \noindent{{\ZSFfontNote #1}}%
  \par\vspace{-\parskip}\vspace{\ZSFspaceS}%
}

% Dünne, graue Trennlinie für Formel-Boxen mit mehreren Elementen
\newcommand{\formulasep}{%
  \par\vspace{-\parskip}\vspace{\ZSFspaceSep}%
  \noindent\textcolor{black!30}{\rule{\linewidth}{0.5pt}}%
  \par\vspace{-\parskip}\vspace{\ZSFspaceSep}%
}


% \runintext: Inline-Text für Tabellen etc. - mit parskip-Kompensation
\newenvironment{runintext}{\par\vspace{-\parskip}\vspace{\ZSFspaceS}\noindent\ignorespaces}{\par}




% ------------------------------------------------------------
% Aussagen (statementbox) und Verfahren (procedure)
% ------------------------------------------------------------
% statementbox: dezente Box mit linkem Farbbalken in Kapitelfarbe.
% Optional fetter Titel in Kapitelfarbe. Fuer Saetze, Definitionen
% und sonstige in sich geschlossene Aussagen, wenn `defbox` mit
% Titelleiste zu schwer wirkt.
\tcbset{zsfstatementstyle/.style={
  enhanced jigsaw,
  breakable=false,
  colback=white,
  colframe=white,
  frame hidden,
  borderline west={\ZSFstatementBarWidth}{0pt}{\StatementBarColor},
  boxrule=0pt,
  arc=0pt, outer arc=0pt,
  left=\ZSFstatementLeftPad,
  right=\ZSFspaceS,
  top=\ZSFspaceXS,
  bottom=\ZSFspaceXS,
  before skip=\ZSFBoxBeforeSkip,
  after skip=\ZSFspaceS,
  before={\ZSFConsumeAfterSubsectionBar},
}}

\NewDocumentEnvironment{statementbox}{O{}}
  {\begin{tcolorbox}[zsfstatementstyle]%
   \if\relax\detokenize{#1}\relax\else
     {\bfseries\color{\StatementTitleColor}#1}\par\vspace{\ZSFspaceXS}%
   \fi\ignorespaces}
  {\end{tcolorbox}}

% procedure: nummerierte Schritt-Liste in einer statementbox.
% Jeder Schritt = fetter Titel + eigene Zeile mit Beschreibung.
% Verwendung: \ProcStep{Titel}{Beschreibung}
\NewDocumentEnvironment{procedure}{O{}}
  {\begin{statementbox}[#1]%
   \begin{enumerate}[leftmargin=1.6em,labelsep=0.4em,itemsep=\ZSFprocItemSep,topsep=0pt,parsep=0pt,partopsep=0pt]}
  {\end{enumerate}\end{statementbox}}

\newcommand{\ProcStep}[2]{\item {\bfseries #1}\\[\ZSFprocStepGap]#2}

% factlist: rahmenlose Liste von Definitions-/Eigenschaftsfakten.
% Farbiger Bullet, fetter Lead-in, em-dash, Beschreibung.
% Verwendung: \ZSFFact{Lead-in}{Beschreibung} - oder rohes \item.
\NewDocumentEnvironment{factlist}{}
  {\begin{itemize}[
     leftmargin=1.2em,
     labelsep=0.4em,
     label={\textcolor{\StatementBarColor}{\textbullet}},
     itemsep=\ZSFprocItemSep,
     topsep=\ZSFspaceXS,
     parsep=0pt,
     partopsep=0pt
   ]}
  {\end{itemize}}

% propertylist: gruppierte Charakterisierungen in einer statementbox.
% Wie factlist, aber zusammengefasst unter optionalem Titel.
\NewDocumentEnvironment{propertylist}{O{}}
  {\begin{statementbox}[#1]\begin{itemize}[
     leftmargin=1.2em,
     labelsep=0.4em,
     itemsep=\ZSFprocItemSep,
     topsep=0pt,
     parsep=0pt,
     partopsep=0pt
   ]}
  {\end{itemize}\end{statementbox}}

\newcommand{\ZSFFact}[2]{\item \textbf{#1}\,---\,#2}

\newtcbox{\ZSFdanger}{colback=red!80!black, colframe=red!80!black, coltext=white, on line, boxsep=\ZSFspaceXS, left=\ZSFspaceS, right=\ZSFspaceS, top=\ZSFspaceXS, bottom=\ZSFspaceXS, boxrule=0pt, arc=1pt, fontupper=\sffamily\bfseries}

% --- Nebenläufiges Bild/Text-Layout (für "Ebene Kurven", Tabellen mit Bildern etc.) ---
\newtcolorbox{splitbox}[1][0.3]{
  blankest,
  sidebyside,
  sidebyside align=top,
  sidebyside gap=\ZSFspaceM,
  lefthand width=#1\linewidth,
  lower separated=false,
  before skip=\ZSFspaceS,
  after skip=\ZSFspaceS,
  boxsep=0pt, left=0pt, right=0pt, top=0pt, bottom=0pt
}

% --- Zentriertes Item-Heading für Listen mit Trennstrich (für Kurven, Profile etc.) ---
\newcommand{\ZSFItemHeading}[1]{%
  \par\vspace{-\parskip}\vspace{\ZSFspaceM}%
  {\centering\textbf{#1}\par}\vspace{-\parskip}\vspace{\ZSFspaceS}%
  \noindent\rule{\linewidth}{0.5pt}\par\vspace{-\parskip}\vspace{\ZSFspaceS}%
}



% Kompakte Grid-Box als tabularray-Container
\newtcolorbox{compactgridbox}[1][]{
  zsfpanelbox,
  colframe=black!20,
  breakable=false,
  title={#1},
  boxsep=0pt, left=0pt, right=0pt, top=\ZSFtextMinPad, bottom=\ZSFtextMinPad,
  zsftitlebar,
  boxrule=0.5pt
}

% Gemeinsamer tabularray-Stil fuer alle Valuegrids
\newcommand{\ZSFvaluegridCommon}{
      colsep=4pt,
      rowsep=3pt,
      hlines,
      vlines,
      rows={valign=m}
}

% Parametrisiertes Valuegrid: #1 = Anzahl Datenspalten, #2 = optionaler Titel
% Erzeugt automatisch `l *{#1}{c}` als colspec
\NewDocumentEnvironment{valuegrid}{m O{}}
  {\begin{compactgridbox}[#2]\begin{tblr}{colspec={l *{#1}{c}},\ZSFvaluegridCommon}}
  {\end{tblr}\end{compactgridbox}}

% Freie colspec-Variante fuer Sonderfaelle
\NewDocumentEnvironment{valuegridcustom}{m O{}}
  {\begin{compactgridbox}[#2]\begin{tblr}{colspec={#1},\ZSFvaluegridCommon}}
  {\end{tblr}\end{compactgridbox}}

% --- Legacy-Aliase (Abwaertskompatibilitaet) ---
\NewDocumentEnvironment{valuegridtwo}{O{}}{\begin{valuegrid}{2}[#1]}{\end{valuegrid}}
\NewDocumentEnvironment{valuegridthree}{O{}}{\begin{valuegrid}{3}[#1]}{\end{valuegrid}}
\NewDocumentEnvironment{valuegridfour}{O{}}{\begin{valuegrid}{4}[#1]}{\end{valuegrid}}
\NewDocumentEnvironment{valuegridfive}{O{}}{\begin{valuegrid}{5}[#1]}{\end{valuegrid}}
\NewDocumentEnvironment{valuegridsix}{O{}}{\begin{valuegrid}{6}[#1]}{\end{valuegrid}}
\NewDocumentEnvironment{valuegridseven}{O{}}{\begin{valuegrid}{7}[#1]}{\end{valuegrid}}


\input{styles/66_index}           % opt-in: Stichwortverzeichnis
\ZSFIndexShowPageNumbertrue       % Abschnitt · (S. Seite)
%
%
%
%
%
%
%
\ifdefined\ZSFShowcaseMode
  \input{styles/65_code_style}
  \input{styles/67_code_comments}
\else
  % \input{styles/65_code_style}    % opt-in: Informatik
  % \input{styles/67_code_comments} % opt-in: annotierte Code-Zeilen
\fi
\input{styles/75_pdf_identity}
% ============================================================
% DOCUMENT SETTINGS
% ============================================================

\hypersetup{colorlinks=true, linkcolor=ChapterColor0, urlcolor=ChapterColor2}

\AtBeginDocument{
  \setlength{\abovedisplayskip}{\ZSFspaceXS}
  \setlength{\belowdisplayskip}{\ZSFspaceXS}
  \setlength{\abovedisplayshortskip}{\ZSFspaceXS}
  \setlength{\belowdisplayshortskip}{\ZSFspaceXS}
  \setlength{\mathsurround}{\ZSFspaceXS}
}

% Gemeinsame Basis fuer alle Listen
\setlist{nosep, itemsep=0pt, parsep=0pt, topsep=\ZSFspaceXS, partopsep=0pt}
% Typspezifische leftmargin (verhindert Duplikation in Kapiteln)
\setlist[itemize]{leftmargin=3mm}
\setlist[enumerate]{leftmargin=4mm}
\setlength{\tabcolsep}{\ZSFspaceM}


% Keine Umbruchsteuerung an Titelbalken in dieser ZSF: weder Platzreserven
% noch der negative Penalty, der vor jedem Balken einen Spaltenumbruch
% anbietet. Die Spalten laufen LaTeX-nativ voll; wo der Umbruch nicht passt,
% steht ein \ZSFNewColumn im Kapitel. Steht hier und nicht bei den Reglern
% oben, weil es ein Befehl ist und 60_boxes dafür geladen sein muss.
%
% Das Umbruchverbot direkt hinter einem Balken bleibt aktiv — ein Balken
% bleibt weiterhin nie als letztes Element einer Spalte stehen.
\ZSFBreakReservesOff
